\section{A045994: a recurrence for a Motzkin convolution}

Let $m_0=m_1=1$ and
\[
 m_{r+2}=m_{r+1}+\sum_{k=0}^{r}m_km_{r-k}\qquad(r\ge0)
\]
be the Motzkin numbers. The defining recurrence of A045994 is
\begin{equation}\label{eq:a045994-source}
 a_0=1,\qquad
 a_n=m_n+\sum_{k=1}^{n-1}m_ka_{n-k-1}\quad(n\ge1).
\end{equation}
In particular, $a_0,\ldots,a_5=1,1,3,7,18,47$.
We prove the recurrence conjectured by R.~J.~Mathar in 2013
\cite{a045994-source}.

\begin{theorem}\label{thm:a045994}
For every integer $n\ge5$, the sequence defined by
\eqref{eq:a045994-source} satisfies
\[
\begin{aligned}
0={}&(n+1)a_n+2(-2n-1)a_{n-1}+2(-n+3)a_{n-2}\\
 &+(11n-19)a_{n-3}+(11n-27)a_{n-4}+3(n-3)a_{n-5}.
\end{aligned}
\]
\end{theorem}

\begin{proof}
Work in the formal power-series ring $\mathbb Q[[x]]$, and put
$M=\sum_{n\ge0}m_nx^n$ and $A=\sum_{n\ge0}a_nx^n$.
The two defining recurrences give
\begin{equation}\label{eq:a045994-gfs}
 M=1+xM+x^2M^2,\qquad A=M+x(M-1)A.
\end{equation}
The second equality follows directly from the Cauchy product:
the coefficient of $x^n$ in $x(M-1)A$ is
$\sum_{k=1}^{n-1}m_ka_{n-k-1}$.
It gives $(1+xA)M=(1+x)A$, or equivalently
$A=M/(1+x-xM)$; the denominator has constant term one.

Define
\[
 p=x^4+3x^3+2x^2-x,\qquad q=x^2+2x-1.
\]
Eliminating $M$ from \eqref{eq:a045994-gfs} yields
\begin{equation}\label{eq:a045994-quadratic}
 pA^2+qA+1=0.
\end{equation}
For an explicit elimination certificate, set $U=1+xA$ and
$V=(1+x)A$. Polynomial expansion gives
\[
\begin{aligned}
pA^2+qA+1={}&
 U^2\{x^2M^2+(x-1)M+1\}\\
 &-(UM-V)\{x^2(UM+V)+(x-1)U\}.
\end{aligned}
\]
Both terms on the right vanish by \eqref{eq:a045994-gfs}.

Differentiating \eqref{eq:a045994-quadratic} gives
\begin{equation}\label{eq:a045994-derivative}
 (2pA+q)A'+p'A^2+q'A=0.
\end{equation}
Now put
\[
\begin{aligned}
u&=x-4x^2-2x^3+11x^4+11x^5+3x^6,\\
v&=1-6x+2x^2+14x^3+17x^4+6x^5,\\
w&=1-4x+x^2,\qquad
L=1-6x+8x^2+2x^3-x^4.
\end{aligned}
\]
A second polynomial certificate is
\[
\begin{aligned}
(2pA+q)(uA'+vA-w)
={}&u\{(2pA+q)A'+p'A^2+q'A\}\\
&+L(pA^2+qA+1).
\end{aligned}
\]
The right side is zero by \eqref{eq:a045994-quadratic} and
\eqref{eq:a045994-derivative}. The constant coefficient of
$2pA+q$ is $-1$, so it is a unit. Cancellation proves the exact
differential equation
\begin{equation}\label{eq:a045994-ode}
 uA'+vA=w.
\end{equation}

For $n\ge6$, coefficient extraction in
\eqref{eq:a045994-ode} gives the asserted recurrence: if
$u=\sum_{j=0}^{5}\alpha_jx^{j+1}$ and
$v=\sum_{j=0}^{5}\gamma_jx^j$, its left coefficient is
$\sum_{j=0}^{5}\{\alpha_j(n-j)+\gamma_j\}a_{n-j}$.
The coefficient of $x^n$ on the right is zero.
The remaining endpoint $n=5$ follows from the six initial values:
\[
6\cdot47-22\cdot18-4\cdot7+36\cdot3+28\cdot1+6\cdot1=0.
\]
This proves every index $n\ge5$.
\end{proof}

\paragraph{Correction to the printed generating function.}
The source also prints $1/(1-x(1+x)M)$, which is incompatible with
its defining recurrence. Since $M=1+x+2x^2+O(x^3)$, that reciprocal
has coefficient $8$ at $x^3$, whereas \eqref{eq:a045994-source}
gives $a_3=7$. The proof above therefore uses the literal source
definition and derives its transform. The discrepancy and the
recurrence are one target family.

\paragraph{Prior work and formal scope.}
The Motzkin generating-function equation is classical
\cite{a045994-motzkin}. The source links Kruchinin and Kruchinin's
general composita methods \cite{a045994-kruchinin}; these are
credited as background. Niu's 2026 treatment of A176677
\cite{a045994-niu} is a related application of algebraic elimination
to a different recurrence. No novelty is claimed for that method.
The bounded source and repository search, including inspected
implementation hits, is recorded in \texttt{source\_checks.md};
it is not a proof of global priority.

The Lean development starts from the two literal natural-number
recurrences, proves \eqref{eq:a045994-gfs}, checks both polynomial
certificates, and extracts coefficients for the full range.
Its final theorems are
\texttt{A045994.source\_recurrence} and
\texttt{A045994.printed\_generating\_function\_false}.
No assertion about the separate path interpretation in the source
comments is included.
